\documentclass[12pt]{article}

\usepackage[utf8]{inputenc}
\usepackage{amsmath, amssymb, amsthm}
\usepackage{geometry}
\geometry{a4paper, margin=2.5cm}

\title{Funcție Aditivă}
\author{}
\date{}

\theoremstyle{definition}
\newtheorem{definition}{Definiție}

\begin{document}

\maketitle

\begin{definition}
Fie $(\Omega, \mathcal{F})$ un spațiu măsurabil. 
O funcție $\mu : \mathcal{F} \to [0, \infty]$ se numește \emph{funcție aditivă} (sau \emph{finit aditivă}) dacă pentru orice două mulțimi disjuncte $A, B \in \mathcal{F}$ avem


\[
    \mu(A \cup B) = \mu(A) + \mu(B).
\]


\end{definition}

\noindent
Echivalent, $\mu$ este aditivă dacă este închisă la reuniuni finite de mulțimi disjuncte:


\[
    \mu\!\left(\bigcup_{i=1}^n A_i\right)
    = \sum_{i=1}^n \mu(A_i),
    \quad A_i \cap A_j = \varnothing \text{ pentru } i \neq j.
\]



\end{document}
